\section{Theorie}
\label{sec:Theorie}


Der Operationsverstärker, kurz OPV, ist ein gleichspannungsgekoppelter Differenzverstärker. Er kann je nach anschluss in vielen bereichen verwendung finden. So kann er als Integriere oder Differenziere 
einer Spannung genutzt werden. Daher Erhält der OPV auch seinen namen. Er kann jedoch auch zur reinen Verstärkung eines Eingangssignals genutzt werden, oder als eine Art schaltvorrichting in form von 
einer Schmidt-Trigger Schaltung. 

\subsection{Idealer und realer Operationsverstärker}
In der Theorie wird oft mit dem Modell des idealen OPV gearbeitet \cite{tietzeschenk}. Dieser zeichnet sich durch eine unendliche Leerlaufverstärkung ($V \to \infty$), 
einen unendlichen eingangswiderstand ($R_e \to \infty$) und einen Ausgangswiderstand von Null ($R_a = 0$) aus.
Reale Operationsverstärker, wie sie in diesem Versuch genutzt werden, Haben aufgrund Technischer limitierungen Grenzen, so Kann nur ein Sehr großer Wiederstand gewählt werden, sodass fast 
kein Strom in den Verstärker fließt. Des weiteren weisen Transistoren physikalische grenzen auf, was zu eier frequenzabhängigen obergrenze für die verstärkug führt. Ab überschreiten dieser Grenze 
können die Transistoren nicht mehr schnell genug schalten und funktionieren nicht mehr ordnungsgemäß. Des weiteren kan es zu einer Offsetspannung kommen, wenn eine Ausgangsspannung anliegt, obwohl es 
keine Differenz im  Eingang gibt.


\subsection{Rückkopplung und Gegenkopplung}
Die Eigenschaften einer OPV-Schaltung hängen maßgeblich von der Art der Rückführung des Ausgangssignals ab \cite{demtroeder2}:
\begin{itemize}
    \item \textbf{Gegenkopplung:} Ein Teil des Ausgangssignals wird auf den invertierenden Eingang ($-$) zurückgeführt. Dies führt zu einer Stabilisierung des Systems, verringert die Gesamtverstärkung und erlaubt den präzisen Betrieb als linearer Verstärker.
    \item \textbf{Mitkopplung:} Die Rückführung erfolgt auf den nicht-invertierenden Eingang ($+$). Dies führt zu einer Instabilität, wodurch der OPV die Differenzspannung maximal verstärkt und nur zwei stabile Zustände (die positive oder negative Sättigung) einnimmt. In dieser Konfiguration lässt sich der OPV als Schalter bzw. Komparator nutzen \cite[siehe auch][]{tietzeschenk}.
\end{itemize}

\subsection{Mathematische Beschreibung verschiedenen Schaltungen}

\subsubsection{Invertierender Linearverstärker}
Durch die hohe Verstärkung und die Gegenkopplung stellt sich am invertierenden Eingang ein virtueller Nullpunkt ein. Das Verhältnis der Widerstände bestimmt den Verstärkungsfaktor
\begin{equation}
    v = \frac{U_a}{U_e} = -\frac{R_2}{R_1}
\end{equation}

\subsubsection{Integrator und Differenzierer}
Wird ein Kondensator mit Kapazität $C$ in die Schaltung integriert, lassen sich zeitabhängige Operationen realisieren. Der Umkehr-Integrator folgt der Beziehung
\begin{equation}
    U_a(t) = -\frac{1}{RC} \int_{0}^{t} U_e(t') \, dt'
\end{equation}
Beim Differenzierer hingegen ist die Ausgangsspannung proportional zur zeitlichen Ableitung der Eingangsspannung.
\begin{equation}
    U_a(t) = -R \cdot C \cdot \frac{d U_e(t)}{dt}
\end{equation}

\subsubsection{Nicht-invertierender Schmitt-Trigger}
Durch die Mitkopplung über einen Spannungsteiler entstehen zwei Schaltschwellen $U_{on}$ und $U_{off}$. Nur bei überwinden der jewailigen Schaltschwelle fällt der OPV in die positive bzw negative Sättignung. Die Differenz dieser Schwellen wird als Hysterese bezeichnet und verhindert ein unerwünschtes Hin- und Herschalten bei verrauschten Signalen.

% --- Beispiel für das Literaturverzeichnis am Ende des Dokuments ---




